% page 210 du fac-simile (image p-209 du scan)
% bloc 157.5 x 250.6 mm ; marge gauche 8.0 mm
\pgc{-12.700mm}{0.000mm}{
\l{\cel{3}2O2}
\l{}
\l{}
\l{\sou{grilio}. (\sou{zool.})\cel{2}Insekto de la klaso \textquotedbl{}ortopteri\textquotedbl{}, che qua la}
\l{\cel{3}maskulo produktas bruiso akuta per interfrotar sua elitri,}
\l{\cel{3}e qua tre prizas la loki suno-radioza, varma. - L.\cel{1}\sou{gryllus}.}
\l{\cel{3}- DFIS.}
\l{}
\l{\sou{griliotalpo}. (\sou{zool.)}\cel{2}Insekto qua devastas la gardeni, vivas}
\l{\cel{3}sub la tersurfaco, e bruisas analoge a grilio. - L.\cel{1}\sou{gryllus}}
\l{\cel{3}\sou{gryllotalpa}.}
\l{}
\l{\sou{grimar}. (\sou{trans.)}\cel{2}(\sou{teatro}).Rugizar (ulu) por igar lu aspektar}
\l{\cel{3}plu olda. - FR.}
\l{}
\l{\sou{grimasar.}\cel{1}(\sou{netr}ans.) Quaze tordar sua vizajo. - DEFR.}
\l{}
\l{grincar. (\sou{netrans}.) (\sou{pri mekanismo)}.\cel{2}Bruisar akre pro tre}
\l{\cel{3}streta kontakto sen interpozeso di oleo-strato. - FI.}
\l{}
\l{\sou{grindar}. (\sou{trans}.)Submisar (korpo generale harda : torn-utensi-}
\l{\cel{3}lo, frezo-cilindro) a la efiko di grosa disko qua rotacas}
\l{\cel{3}rapide (ordinare ek petro tre harda : karborundo, smerilo,}
\l{\cel{3}diamanto-polvo) por quaze limar lu.}
\l{}
\l{\sou{gripo}. (\sou{patol.}) Quaza kataro epidemiala. - DEFIRS.}
\l{}
\l{\sou{griza}.\cel{2}Koloro meza inter blanka e nigra. - EFIS.}
\l{}
\l{\sou{grogo}.\cel{2}Drinkajo ek aquo kolda o varma, en qua onu varsas su-}
\l{\cel{3}kro, brandio od altra alkoholajo, adjuntante citrono-}
\l{\cel{3}lonchi. - DEFIR.}
\l{}
\l{\sou{gronda}r. (\sou{netr}ans.) Audigar bruiso duroza e nelauta, quale}
\l{\cel{3}la tondro fora od ula animali iracoza. - DEFIRS.}
\l{}
\l{\sou{gropo}.\cel{2}(\sou{anat.)}\cel{3}Che ula animali, parto dopa, sferatra, de}
\l{\cel{3}la hancho til la loko ube komencas la kaudo. - DEFIRS.}
\l{}
\l{\sou{grosa}.\cel{2}I. Qua transiras la volumino ordinara. - II. Qua}
\l{\cel{3}transiras la volumino di altra kozo sama-speca. - EFIS.}
\l{}
\l{\sou{grosdekeduo}.\cel{2}Dekeduo de dekedui. (t. e. l44).}
\l{}
\l{\sou{grosiera}. I. (\sou{pri kozo})\cel{2}Facita ye maniero vulgara, o fasonita}
\l{\cel{3}krude. - II. (\sou{pri person}o)Di qua la rudeso ne pludolcigesas}
\l{\cel{3}da civilizeso. - EFIS.}
\l{}
\l{\sou{groto}. Exkavajo pitoreska, naturala o manuo-facita. - DEFIRS.}
\l{}
\l{\sou{groteska}.\cel{2}Rid-ecitanta pro lua aspekto baroka, stranja.}
\l{\cel{3}- DEFIRS.}
\l{}
\l{\sou{grozelo}. (\sou{bot.})\cel{2}Frukto redatra o verda, plu grosa kam la ribo}
\l{\cel{3}(uzata, kande lu esas ankore nematura ed acido-saporanta,}
\l{\cel{3}por facar sauco qua akompanas makrelo-disho) qua kreskas}
\l{\cel{3}sur arbusto qua formacas genero de la familio \textquotedbl{}grosulariei\textquotedbl{}}
\l{\cel{3}- L.\cel{1}\sou{Ribes grossularia.}\cel{1}- FS.}
}
